% Part: intuitionistic-logic
% Chapter: tableaux
% Section: introduction

\documentclass[../../../include/open-logic-section]{subfiles}

\begin{document}

\olfileid{int}{tab}{int}

\olsection{سريزه}

!!^{tableau}s د !!{signed
  formula}s ځانګړې، لاندې خوا ته څانګېدونکې ونې دي۔ هره نښه‌لرونکې
جوړه د صدق نښه ($\True$ يا $\False$) او !!a{sentence} لري:
\[
\sFmla{\True}{!A} \text{ يا } \sFmla{\False}{!A}.
\]
!!^a{tableau} له څو \emph{فرضونو} پيلېږي۔ هر بل
!!{signed formula} د استنتاج د يوې قاعدې په تطبيقولو جوړېږي۔
ځينې قاعدې د ونې يوه څوکه ته يو يا څو !!{signed formula}s
ورزياتوي؛ نورې دوه نوې څوکې ورزياتوي او دوه څانګې جوړوي۔ د
قاعدې پايله داسې !!{signed formula}s وي چې !!{formula} يې تر
هغه !!{signed formula} ساده وي چې قاعده پرې تطبيق شوې وي۔
کله چې په يوه څانګه کښې $\sFmla{\True}{!A}$ او
$\sFmla{\False}{!A}$ دواړه وي، هغه څانګه \emph{تړلې} بولو۔
که د !!a{tableau} هره څانګه تړلې وي، ټول !!{tableau} تړلے
دے۔ تړلے !!{tableau} داسې !!a{derivation} جوړوي چې ښيي د
!!{tableau} د پيل !!{signed formula}s سټ نۀ شي پوره کېدے۔ له دې څخه د
$\Proves$ اړيکه تعريفولے شو: $\Gamma \Proves !A$ هله او
يوازې هله چې داسې متناهي سټ~$\Gamma_0 = \{!B_1,
\dots, !B_n\} \subseteq \Gamma$ شته چې د لاندې فرضونو دپاره
تړلے !!{tableau} وي:
\[
\{\sFmla{\False}{!A}, \sFmla{\True}{!B_1}, \dots, \sFmla{\True}{!B_n}\}.
\]

د شهودي منطق دپاره بايد د !!{signed formula} مفهوم هم پراخ
کړو او د نښلوونکو قاعدې هم بدلې کړو۔ د نښې ($\True$ يا
$\False$) تر څنګ، په شهودي !!{tableau}s کښې
!!{formula}s \emph{مخکښونه}~$\sigma$ هم لري۔ د ژباړې د
نسخې څرګندونه (OLINT-012): منجمده سرچينه دلته «مودالي»
تابلوګانې يادوي؛ د دې څپرکي او ورپسې شهودي قاعدو له مخې
«شهودي» مراد دي۔ مخکښونه د مثبتو صحيح عددونو ناتشې
لړۍ دي، يعنې $\sigma \in
(\PosInt)^* \setminus \{\emptyseq\}$۔ دا نښې د چاپ پر مهال له
چاپېرې $\tuple{\ }$ پرته ليکو او بېل !!{element}s يې د
$.$ په وسيله جلا کوو، نه د $,$ په وسيله۔ که $\sigma$ يوه
مخکښ وي، نو $\sigma.n$ د $\sigma \concat \tuple{n}$
معنا لري؛ مثلاً که $\sigma = 1.2.1$ وي، نو $\sigma.3$
به $1.2.1.3$ وي۔ د بېلګې په توګه،
\[
\sFmla{\True}{!A \lif (!B \lif !C)}[1.2]
\]
يو \emph{مخکښ‌لرونکے !!{signed formula}} دے (يا په لنډ ډول
يو \emph{مخکښ‌لرونکے !!{formula}})۔

په تصور کښې مخکښ د مدل يوه نړۍ نوموي چې ښايي د
!!a{tableau} د يوې څانګې !!{formula}s پوره کړي؛ که
$\sigma$ کومه نړۍ ونوموي، نو $\sigma.n$ له~$\sigma$ نومولې نړۍ څخه
لاس‌رسي وړ نړۍ نوموي۔

په شهودي مدلونو کښې د لاسرسي اړيکه انعکاسي او متعديه ده۔ د
د مخکښونو په ژبه، له~$\sigma$ څخه $\sigma$ پخپله لاس‌رسي وړ ده،
او هره هغه نښه هم لاس‌رسي وړ ده چې~$\sigma$ غځوي، يعنې
هره د $\sigma.n_1.\cdots.n_k$ په بڼه نښه۔ د~$\sigma$ خپل
ځان او د هغې هر غځول د ښودلو دپاره $\sigma.*$ نښه کاروو۔ په
بله وينا، د $\sigma.*$ مخکښونه ټول او يوازې هغه مخکښونه دي
چې له~$\sigma$ څخه لاس‌رسي وړ دي۔

\end{document}
